\par\noindent\begin{minipage}{\linewidth}
% !TeX root = ../教材题库_全部条目_转写工作稿.tex
% Source: XJ-B1-C01-TH-P016-Y0927
% 第1章 / 1.2 常用逻辑用语 / 教材 p.16 / PDF 第23页
% 已对照原页转写题干及公式；待教师复核；未制作答案
% 原题及相邻原解截图（核对用）：../题目原图/XJ-B1-C01-TH-P016-Y0927-01.webp, ../题目原图/XJ-B1-C01-TH-P016-Y0927-02.webp
% 本文件只在转写完成后作为题干使用；OCR 条目尚不是可靠题干。
\begin{bankitem}{p.16 思考}{第1章 / 1.2 常用逻辑用语 / 教材 p.16 / PDF 第23页}{XJ-B1-C01-TH-P016-Y0927}{已对照原页转写题干及公式；待教师复核；未制作答案}
对于 1.2.1 小节例 2 后面列出的 6 个命题，哪些是真命题？哪些是假命题？
\par\textit{所引用的教材 p.15 六个命题如下：}
\subq{1}若两个三角形全等，则它们相似；
\subq{2}若两个三角形相似，则它们全等；
\subq{3}若实数 $a\ne0$，则 $a^2>0$；
\subq{4}若四边形 $ABCD$ 为菱形，则 $AC\perp BD$；
\subq{5}若 $a<0$，则方程 $ax^2+4x-3=0$ 没有正的实根；
\subq{6}若 $a=b$，则 $a^2=b^2$。
\end{bankitem}
\end{minipage}\par\medskip
